\documentclass[11pt,a4paper]{article}
\usepackage[utf8]{inputenc}
\usepackage[T1]{fontenc}
\usepackage[portuguese]{babel}
\usepackage{xcolor}
\usepackage{hyperref}
\usepackage{geometry}
\usepackage{listings}
\usepackage{tcolorbox}
\geometry{margin=2.5cm}
\definecolor{primary}{RGB}{41,128,185}
\definecolor{secondary}{RGB}{44,62,80}
\definecolor{codebg}{RGB}{240,240,240}
\lstset{backgroundcolor=\color{codebg},basicstyle=\ttfamily\small,breaklines=true,frame=single}
\title{\color{secondary}\Huge\textbf{Solução: Exercício 2.5 - StatefulSet e Headless Service}}
\author{\color{primary}DevOps com Kubernetes -- 2026}
\date{}
\begin{document}
\maketitle


\textbf{Guia de Solução:}

\subsection*{\color{secondary}Passo a Passo}

1. Crie um Headless Service com \texttt{clusterIP: None}.
2. Crie um StatefulSet com \texttt{serviceName} apontando para esse Service.
3. Use duas réplicas e confirme os nomes \texttt{app-0} e \texttt{app-1}.
4. Teste resolução DNS com \texttt{nslookup app-0.nome-do-service}.

\end{document}
